\section{Limitations, Ethics, and Conclusion}
\label{sec:limitations}

\subsection{Limitations}

The common evaluation uses one recorded generation workflow and one fully evaluated object class; exact editor provenance for the 128 restaurant pairs requires manifest reconciliation. It does not yet test laundering, a held-out editor, whole-image contrast sets, real-object paste-back, genuine-object controls, or humans. Five pristine sources repeat, and source-cluster uncertainty remains to be computed. The selected cat and trash-can expansions are excluded pending common canonicalization and resolution of a cross-decoder pixel-integrity audit. The legacy MLLM export is format-confounded, and commercial results are heterogeneous diagnostics rather than a balanced leaderboard. Accordingly, our conclusion is limited to the canonical mouse condition.

\subsection{Ethics and Release}

Plausible manipulated evidence is dual-use. CLAIMFORGE is intended to help platforms, insurers, and forensic researchers evaluate defenses, but unrestricted release could also lower the cost of misuse. We therefore plan a gated, research-use release with source-license documentation, provenance, and explicit restrictions on operational fraud. This policy remains a release plan until implemented.

\subsection{Conclusion}

Relative to the decoded pre-canonical source, a median manipulated pair remains 99.887\% unchanged while its apparent claim-relevant meaning changes. On 275 paired mouse edits, current methods expose useful but incomplete forensic signal: the strongest joint model misses the image-level reference, and the strongest continuous localizer lacks a registered detector. The immediate next steps are source-cluster uncertainty, provenance reconciliation, audited multi-object canonicalization, and realistic laundering tests.
